\chapter{Den förenklade kontrollern}
\label{ch:kursen}

{\large\itshape Vad en undervisningskontroller implementerar, och vad den medvetet låter bli.\par}

\begin{chapterbox}{I det här kapitlet}
\begin{itemize}
  \item Blocken en förenklad CAN-kontroller består av.
  \item Vilka delar av \kapref{ch:buss}--\kapref{ch:fel} som finns med, och vilka som inte gör det.
  \item Varför en pollande mjukvara inte kan prata direkt med kontrollern.
  \item Registerlagret som gör den nåbar, och de tre begränsningar som når ända upp i koden.
\end{itemize}
\end{chapterbox}

De sex föregående kapitlen beskriver CAN som standarden definierar det. Det här kapitlet
beskriver den kontroller som faktiskt byggs i undervisningen: vad den gör av allt det, och vad
den låter bli.

Kapitlet vänder sig åt båda hållen. Bygger man kontrollern är listan nedan kravspecifikationen.
Skriver man drivrutinen är den i stället listan över vad man inte får räkna med.

\section{Blocken}

\bookfigure{controller_blocks}{Den förenklade kontrollerns block, och lagret som gör den nåbar
från en processor.}{fig:blocks}

\begin{booktable}{@{}lL@{}}
\thead{Block} & \thead{Ansvar} \\ \midrule
Bittimer & Räknar tick per bit, pulsar vid sampelpunkten och vid bitperiodens slut
  (\kapref{ch:tajming}). \\
CRC-15 & En bitseriell motor som både beräknar och kontrollerar (\kapref{ch:stuffing}). \\
Sändningsskiftregister & Serialiserar MSB först och sätter in stoppbitar. \\
Mottagningsskiftregister & Deserialiserar och plockar bort stoppbitar; rapporterar stoppfel. \\
Kontroller & Tillståndsmaskinen som sekvenserar framens fält och sköter arbitreringen. \\
Synkroniserare & Två vippor på varje asynkron ingång, så att bussen och reset kan läsas säkert. \\
\end{booktable}

Ovanpå dem ligger två block till, som inte har med CAN att göra utan med att göra kontrollern
\emph{nåbar}:

\begin{booktable}{@{}lL@{}}
\thead{Block} & \thead{Ansvar} \\ \midrule
Registerbank & Översätter kontrollerns encykelspulser till klibbiga, pollbara register. \\
Serie\-brygga & Gör registren läs- och skrivbara utifrån, över en seriebuss. \\
\end{booktable}

\section{Vad som finns med}

Kontrollern implementerar framens sekvens i sin helhet: SOF, arbitreringsfältet med identifierare
och RTR, kontrollfältet med IDE, r0 och DLC, datafältet, CRC-15 med avgränsare, ACK-lucka med
avgränsare, och sju bitars EOF.

Den implementerar också:

\begin{itemize}
  \item \textbf{Bitstoppning åt båda hållen}, med stoppningstillståndet levande över fältgränser.
  \item \textbf{CRC-15}, med samma motor på sändar- och mottagarsidan.
  \item \textbf{Arbitrering}, bit för bit över arbitreringsfältet.
  \item \textbf{Öppen dränering}: noden driver bara när den ska, och bussen är dominant om någon
    drar ned den.
  \item \textbf{Stoppfelsdetektering} och \textbf{CRC-kontroll} på mottagarsidan.
\end{itemize}

Det är tillräckligt för att två sådana noder ska utbyta riktiga CAN-frames på en riktig buss, och
tillräckligt för att en kommersiell bussanalysator ska kunna avkoda dem.

\section{Vad som inte finns med}

\begin{booktable}{@{}lL@{}}
\thead{Saknas} & \thead{Konsekvens} \\ \midrule
Felframes & Ett fel avbryter inte överföringen för de andra noderna. \\
Felräknare & Inget error passive, inget bus off (\kapref{ch:fel}). \\
Automatisk omsändning & En frame som gick fel skickas inte om av hårdvaran. \\
ACK-återläsning & Sändaren får aldrig veta om någon kvitterade. \\
Mottagning under sändning & En nod som förlorar arbitreringen tar inte emot framen som vann. \\
Mottagningskö & Kontrollern håller exakt en mottagen frame. \\
Filtrering & Varje mottagen frame lämnas vidare; filtreringen ligger i mjukvaran. \\
Utökade identifierare & Bara 11 bitar. \\
Bitfelsdetektering utanför arbitreringen & Sändaren läser bara tillbaka bussen under arbitreringsfältet. \\
\end{booktable}

Var och en av dem är ett medvetet val. Tillsammans är de skillnaden mellan en konstruktion som
går att bygga och verifiera under en kurs, och en som inte gör det.

\section{Varför ett registerlager behövs}

Kontrollern är byggd för en anropare i samma klockdomän. Den signalerar med pulser som är
\textbf{en klockcykel} långa --- 20~ns vid 50~MHz.

En processor som pollar över en seriebuss ser systemet några tiotals mikrosekunder i taget i
bästa fall. En puls på 20~ns är alltså borta tusentals cykler innan nästa pollning kommer.

\textbf{Registerbanken löser det genom att göra om pulser till nivåer.} En puls sätter en bit,
och biten står kvar tills mjukvaran uttryckligen nollställer den:

\begin{booktable}{@{}lll@{}}
\thead{Statusbit} & \thead{Sätts av} & \thead{Nollställs av} \\ \midrule
TX klar & en avslutad sändning & en accepterad sändningstrigger \\
RX giltig & en mottagen frame & en kvittering från mjukvaran \\
Fel & en stigande flank på felnivån & en uttrycklig nollställning \\
\end{booktable}

Det mönstret --- \emph{klibbig bit, uttrycklig nollställning} --- är hela poängen med lagret, och det
är också den vanligaste källan till buggar i mjukvaran ovanför. Glöms nollställningen står biten
kvar för alltid, och drivrutinen läser samma frame i en oändlig loop.

\begin{aside}
\note Registerbanken \emph{döljer} inte kontrollerns begränsningar, den \emph{vidarebefordrar} dem.
Det är ett designval värt att lägga märke till: ett lager som hade buffrat mottagna frames hade
sett hjälpsamt ut och gjort systemets verkliga beteende omöjligt att resonera om.
\end{aside}

\section{De tre begränsningarna som når ända upp}

Tre av förenklingarna ovan är inte bara interna. De syns i mjukvaran, och den som skriver
drivrutinen måste hantera dem snarare än att bli förvånad av dem.

\subsection{En förlorad arbitrering ser ut som en lyckad sändning}

Kontrollern signalerar "sändningen är över" i båda fallen. Biten betyder \textbf{porten är fri
igen}, inte \emph{framen kom fram}.

Att den gör så är avsiktligt, och det är faktiskt en rättelse av en verklig bugg: signalerade den
bara vid lyckad sändning, skulle biten bli liggande efter varje förlorad arbitrering, och noden
kunde inte sända igen förrän den nollställdes för hand. På en tvånodsbuss alltså i normal drift.

Konsekvensen för mjukvaran är en pollningsloop i två steg:

\begin{console}
vänta på "porten är fri"   -> sändningsförsöket är över
läs felflaggan             -> gick det bra eller inte?
\end{console}

Utan det andra steget går det inte att veta om framen ska skickas om.

\subsection{En bit för fyra orsaker}

Förlorad arbitrering, stoppfel, misslyckad CRC och mottagen fjärrframe ger alla samma enda
felflagga. Mjukvaran kan se \emph{att} något gick fel, aldrig \emph{vad}.

Retry-logik som behandlar en förlorad arbitrering annorlunda än en trasig CRC går alltså inte att
skriva mot det här gränssnittet. Rätt hantering är att veta om begränsningen och behandla alla
fyra lika --- inte att gissa orsaken ur sammanhanget.

\subsection{Ingen buffring på mottagarsidan}

Kontrollern håller en frame. Kommer nästa innan den förra kvitterats skrivs den över, och den
nyaste vinner.

Det är dokumenterat och avsiktligt. En riktig kontroller har en mottagningskö; den här har det
inte, och mjukvaran får polla tillräckligt ofta. Vid någon trafikintensitet räcker det inte
oavsett hur ofta den pollar, och då tappas frames --- vilket är hårdvarans begränsning och inte
mjukvarans bugg.

\section{Vad en riktig kontroller skulle lägga till}

För fullständighetens skull, och som svar på frågan "hur långt är det kvar?":

\begin{itemize}
  \item Felframes, felräknare och bus off (\kapref{ch:fel}).
  \item Automatisk omsändning efter fel.
  \item Mottagningsfilter i hårdvara, så att ointressanta frames aldrig når mjukvaran.
  \item En mottagningskö med flera platser.
  \item Flera sändningsbuffertar med inbördes prioritet.
  \item Utökade identifierare, och därmed CAN 2.0B.
  \item Återläsning av hela framen, och därmed riktig bitfelsdetektering.
\end{itemize}

Listan är lång, men lägg märke till vad som \emph{inte} står på den: ramning, stoppning, CRC,
arbitrering och bittajming. Det svåra i CAN är redan gjort.

\section*{Övningar}

\exercise{Vad saknas}{Resonemang}
\label{ex:7:1}
För vart och ett av följande scenarier: vad gör en fullständig CAN-kontroller, och vad gör den
förenklade?
\begin{enumerate}[a)]
  \item En frame förlorar arbitreringen.
  \item En mottagen frame har fel checksumma.
  \item Ingen nod kvitterar en sänd frame.
  \item Två frames tas emot tätt efter varandra.
\end{enumerate}

\exercise{Klibbiga bitar}{Resonemang}
\label{ex:7:2}
Kontrollern pulsar i en klockcykel; mjukvaran pollar var fyrtionde mikrosekund.
\begin{enumerate}[a)]
  \item Hur många klockcykler på 50~MHz går det mellan två pollningar?
  \item Vad skulle mjukvaran se om statusbiten speglade pulsen direkt?
  \item Varför måste biten nollställas av en skrivning i stället för automatiskt vid läsning?
\end{enumerate}

\exercise{Två steg}{Resonemang}
\label{ex:7:3}
En drivrutin väntar på att sändningen ska bli klar och returnerar sedan "det gick bra".
\begin{enumerate}[a)]
  \item Vilket fel gör den?
  \item Vad måste den göra i stället?
  \item Vilka fyra orsaker kan ligga bakom en satt felflagga?
\end{enumerate}

\exercise{Gränsen}{Resonemang}
\label{ex:7:4}
Registerbanken vidarebefordrar kontrollerns begränsningar i stället för att dölja dem.
\begin{enumerate}[a)]
  \item Ge ett exempel på hur den \emph{skulle} kunna dölja en av dem.
  \item Vad hade det kostat den som felsöker systemet?
  \item Vilken princip säger detta något om, mer allmänt, när man konstruerar ett gränssnitt
    mellan hårdvara och mjukvara?
\end{enumerate}
